\documentclass[12pt,a4paper]{article}
\usepackage[margin=25mm, headheight=15pt, headsep=10pt]{geometry}
\usepackage{fontspec}
\usepackage[table]{xcolor} % Note: table option is required for rowcolors
\usepackage{titlesec}
\usepackage{tocloft}
\usepackage{fancyhdr}
\usepackage{booktabs}
\usepackage{tcolorbox}
\tcbuselibrary{skins, breakable}
\usepackage[colorlinks=true, linkcolor=emerald, urlcolor=emerald, bookmarks=true]{hyperref}

\definecolor{emerald}{RGB}{0,102,51}
\definecolor{darkred}{RGB}{139,0,0}
\definecolor{lightgray}{RGB}{245,245,245}

\usepackage[nil]{babel}
\babelprovide[import, main]{bengali}
\babelprovide[import]{arabic}
\babelfont{rm}[Renderer=HarfBuzz,Scale=1.0]{Noto Serif Bengali}
\babelfont{sf}[Renderer=HarfBuzz]{Noto Sans Bengali}
\babelfont{tt}[Renderer=HarfBuzz]{Noto Sans Bengali}
\babelfont[arabic]{rm}[Renderer=HarfBuzz,Scale=1.1]{Noto Naskh Arabic}

% Section styling
\titleformat{\section}{\Large\bfseries\color{emerald}}{\thesection.}{0.5em}{}
\titleformat{\subsection}{\large\bfseries\color{emerald!85!black}}{\thesubsection.}{0.5em}{}
\titleformat{\subsubsection}{\normalsize\bfseries\color{emerald!70!black}}{\thesubsubsection.}{0.5em}{}
\titlespacing*{\section}{0pt}{18pt}{8pt}

% Fancy Header/Footer setup
\pagestyle{fancy}
\fancyhf{} % clear all header and footer fields
\fancyhead[L]{\color{emerald}\small\textbf{\nouppercase{\leftmark}}}
\fancyfoot[L]{\scriptsize\color{black!50}{\lat{AI}}-সংকলিত, আলেম-যাচাইকৃত নয়}
\fancyfoot[C]{\color{emerald!70!black}\thepage}
\renewcommand{\headrulewidth}{0.4pt}
\renewcommand{\headrule}{\hbox to\headwidth{\color{emerald}\leaders\hrule height \headrulewidth\hfill}}
\renewcommand{\footrulewidth}{0pt}

% Custom boxes
% 1. Ayat/Hadith Box
\newtcolorbox{ayatbox}[1][]{
    enhanced, breakable, 
    colback=emerald!5, 
    colframe=emerald, 
    boxrule=0.5pt, 
    arc=3pt, 
    left=10pt, right=10pt, top=10pt, bottom=10pt,
    #1
}

% 2. Quote/Translation Box
\newtcolorbox{quotebox}[1][]{
    enhanced, breakable,
    frame hidden,
    colback=lightgray,
    borderline west={3pt}{0pt}{emerald},
    left=15pt, right=10pt, top=10pt, bottom=10pt,
    sharp corners,
    #1
}

% 3. AI-generated notice box
\newtcolorbox{warnbox}[1][]{
    enhanced, breakable,
    colback=red!4,
    colframe=darkred,
    boxrule=0.8pt,
    arc=3pt,
    left=10pt, right=10pt, top=10pt, bottom=10pt,
    #1
}

% Commands
\newcommand{\arab}[1]{\foreignlanguage{arabic}{#1}}
\newcommand{\kib}[1]{\textcolor{emerald}{\textbf{#1}}}

% Latin (English) fallback: Noto Serif Bengali has NO Latin glyphs
\newfontfamily\latfont[Renderer=HarfBuzz]{Noto Serif}
\newcommand{\lat}[1]{{\latfont #1}}

\setlength{\parskip}{5pt}
\setlength{\parindent}{0pt}

% Fix for missing bullet in Noto Serif Bengali
\renewcommand{\labelitemi}{$\bullet$}



\begin{document}
\begin{center}{\Large\bfseries\color{emerald} \lat{12}-শেষ-দোয়া-সালাম}\end{center}
\vspace{1em}
\section*{১২ — শেষ দোয়া ও সালাম}

\begin{warnbox}
 \textbf{গুরুত্বপূর্ণ নোটিশ:} এই কন্টেন্টটি \lat{AI} দিয়ে প্রস্তুত। কেবল সহীহ/হাসান হাদিস রাখা হয়েছে।
\end{warnbox}

\medskip\hrule\medskip

\subsection*{১. সূত্র কার্ড}

\begin{table}[h]\centering\small
\begin{tabular}{ll}
বিষয় & বিবরণ \\
\hline
\textbf{বিষয়} & নামাজের শেষে দোয়া ও সালাম \\
\textbf{সূত্র} & সহীহ মুসলিম ৫৮২, সহীহ বুখারি ৮২৯ \\
\textbf{মর্যাদা} & \textbf{সুন্নাতে মুআক্কাদাহ} — সালাম ছাড়া নামাজ অসম্পূর্ণ \\
\hline
\end{tabular}
\end{table}

\medskip\hrule\medskip

\subsection*{২. মূল আরবি}

\textbf{শেষ দোয়া:} \arab{اللَّهُمَّ} \arab{أَعِنِّي} \arab{عَلَىٰ} \arab{ذِكْرِكَ} \arab{وَشُكْرِكَ} \arab{وَحُسْنِ} \arab{عِبَادَتِكَ}

\textbf{সালাম:} \arab{السَّلَامُ} \arab{عَلَيْكُمْ} \arab{وَرَحْمَةُ} \arab{اللَّهِ}

\medskip\hrule\medskip

\subsection*{৩. বাংলা উচ্চারণ}

\textbf{শেষ দোয়া:} আল্লাহুম্মা আ'ইন্নী আলা যিকরিকা ওয়া শুকরিকা ওয়া হুসনি ইবাদাতিক

\textbf{সালাম:} আসসালামু আলাইকুম ওয়া রাহমাতুল্লাহ

\medskip\hrule\medskip

\subsection*{৪. শব্দ-শব্দ অর্থ}

\begin{table}[h]\centering\small
\begin{tabular}{lll}
আরবি & বাংলা & \lat{English} \\
\hline
\textbf{\arab{أَعِنِّي}} & সাহায্য কর & \textit{\lat{help} \lat{me}} \\
\textbf{\arab{ذِكْرِكَ}} & তোমার যিকর & \textit{\lat{remembering} \lat{You}} \\
\textbf{\arab{شُكْرِكَ}} & তোমার শুকর & \textit{\lat{thanking} \lat{You}} \\
\textbf{\arab{حُسْنِ} \arab{عِبَادَتِكَ}} & সুন্দর ইবাদত & \textit{\lat{excellent} \lat{worship}} \\
\textbf{\arab{السَّلَامُ} \arab{عَلَيْكُمْ}} & তোমাদের ওপর শান্তি & \textit{\lat{Peace} \lat{be} \lat{upon} \lat{you}} \\
\textbf{\arab{وَرَحْمَةُ} \arab{اللَّهِ}} & ও আল্লাহর দয়া & \textit{\lat{and} \lat{Allah's} \lat{mercy}} \\
\hline
\end{tabular}
\end{table}

\medskip\hrule\medskip

\subsection*{৫. পূর্ণ বাংলা অনুবাদ}

\textbf{\lat{“}হে আল্লাহ! তোমার যিকর, তোমার শুকর ও তোমার সুন্দর ইবাদতে আমাকে সাহায্য কর।\lat{”}}

\textbf{\lat{“}তোমাদের ওপর আল্লাহর শান্তি ও দয়া বর্ষিত হোক।\lat{”}}

\medskip\hrule\medskip

\subsection*{৬. সংক্ষিপ্ত ব্যাখ্যা}

\subsubsection*{\textbf{সালাম কী?}}
\textbf{\lat{Salam} (সালাম)} মানে শান্তি। বাম দিকে তাকিয়ে "আসসালামু আলাইকুম ওয়া রাহমাতুল্লাহ" বলুন, এরপর ডান দিকে। সালাম দিলে নামাজ \textbf{সম্পূর্ণ} হয় — এটি নামাজের \textbf{\lat{wajib} (ওয়াজিব)} অংশ।

\subsubsection*{\textbf{সালাম কাকে দিচ্ছেন?}}
ফেরেশতা ও উপস্থিত সকল মুসলিমের প্রতি — সেই \textbf{\lat{peace} (শান্তি)} বিনিময় যা জান্নাতের প্রতীক।

\subsubsection*{\textbf{কীভাবে সালাম দিবেন?}}
\begin{itemize}
\item \textbf{ডান দিকে} মাথা ঘোরান — ঠোঁট কাঁধের সমান (অতিরিক্ত নয়)।
\item \textbf{বাম দিকে} মাথা ঘোরান — চিবুকের দিকে।
\item \textbf{ধীরে, স্পষ্ট করে} — এবং মনে করুন: "আমি সকল মুসলিমকে শান্তি দিচ্ছি।"
\end{itemize}
\medskip\hrule\medskip

\subsection*{৭. সংশ্লিষ্ট হাদিস}

\begin{table}[h]\centering\small
\begin{tabular}{ll}
হাদিস & গ্রন্থ \\
\hline
\textbf{\lat{“}নামাজের শেষে ডান ও বাম দিকে সালাম দাও।\lat{”}} & সহীহ মুসলিম ৫৮২ \\
\textbf{\lat{“}সালাম ছাড়া নামাজ পূর্ণ হয় না।\lat{”}} & সুনানে আবু দাউদ ১০১৮ (সহীহ) \\
\hline
\end{tabular}
\end{table}

\medskip\hrule\medskip

\textit{শেষ আপডেট: আগস্ট ২০২৬}

\end{document}
